%==============================================================================
%  Materi Pembelajaran MK Cyber Security
%  Review Paper State-of-the-Art: Federated Learning untuk NIDS
%  Paper yang diulas: Fed-ANIDS (Expert Systems with Applications, 2023)
%==============================================================================
\documentclass[11pt,a4paper]{article}

\usepackage[utf8]{inputenc}
\usepackage[T1]{fontenc}
\usepackage[bahasa]{babel}
\usepackage{amsmath,amssymb}
\usepackage{geometry}
\geometry{margin=2.5cm}
\usepackage{booktabs}
\usepackage{array}
\usepackage{enumitem}
\usepackage{xcolor}
\usepackage{listings}
\usepackage{tikz}
\usetikzlibrary{arrows.meta,positioning,calc,fit,backgrounds,shapes.geometric}
\usepackage{hyperref}
\hypersetup{colorlinks=true,linkcolor=blue,urlcolor=blue,citecolor=blue}

% ----- Gaya node diagram -----
\tikzset{
  server/.style={rectangle,rounded corners,draw=black,fill=blue!12,
                 minimum width=3.3cm,minimum height=1.0cm,align=center,font=\small},
  client/.style={rectangle,rounded corners,draw=black,fill=green!12,
                 minimum width=2.6cm,minimum height=1.0cm,align=center,font=\small},
  store/.style={cylinder,shape border rotate=90,draw=black,fill=orange!15,
                aspect=0.3,minimum width=2.2cm,minimum height=1.1cm,align=center,font=\footnotesize},
  block/.style={rectangle,draw=black,fill=gray!8,minimum width=1.9cm,
                minimum height=0.9cm,align=center,font=\footnotesize},
  lat/.style={circle,draw=black,fill=purple!15,minimum size=0.9cm,font=\footnotesize},
  flow/.style={-{Stealth[length=2.5mm]},thick},
  flowb/.style={{Stealth[length=2.5mm]}-{Stealth[length=2.5mm]},thick},
}

% ----- Gaya listing kode Python -----
\definecolor{codegray}{gray}{0.95}
\definecolor{codegreen}{rgb}{0.0,0.5,0.0}
\definecolor{codeblue}{rgb}{0.0,0.0,0.7}
\definecolor{codepurple}{rgb}{0.58,0.0,0.82}
\lstdefinestyle{py}{
  language=Python,
  backgroundcolor=\color{codegray},
  basicstyle=\ttfamily\footnotesize,
  keywordstyle=\color{codeblue}\bfseries,
  stringstyle=\color{codegreen},
  commentstyle=\color{codepurple}\itshape,
  numberstyle=\tiny\color{gray},
  numbers=left,
  numbersep=6pt,
  frame=single,
  breaklines=true,
  showstringspaces=false,
  tabsize=2,
  captionpos=b
}

\title{\textbf{Review Paper State-of-the-Art:\\
Federated Learning untuk Network Intrusion Detection System (NIDS)}\\[4pt]
\large Studi Kasus Metode Unggulan: \textit{Fed-ANIDS}}
\author{\textit{Cyber Security}}
\date{}

\begin{document}
\maketitle

%------------------------------------------------------------------
\begin{center}
\begin{tabular}{@{}ll@{}}
\textbf{Nama} & : Hero Yudo Martono \\
\textbf{NRP}  & : 1026900012 \\
\end{tabular}
\end{center}
\vspace{0.5em}
\hrule
\vspace{1em}

%==================================================================
\section{Identitas Penulis (Paper yang Diulas)}
%==================================================================

Paper yang dipilih sebagai representasi metode unggulan FL untuk NIDS adalah
sebagai berikut.

\begin{table}[h]
\centering
\renewcommand{\arraystretch}{1.35}
\begin{tabular}{@{}>{\raggedright\arraybackslash}p{4.2cm}>{\raggedright\arraybackslash}p{9.5cm}@{}}
\toprule
\textbf{Komponen} & \textbf{Keterangan} \\
\midrule
Judul makalah &
\textit{Fed-ANIDS: Federated learning for anomaly-based network intrusion
detection systems} \\
Penulis &
Meryem Janati Idrissi, Hamza Alami, Abdelkader El Mahdaouy, Abdellah El Mekki,
Soufiane Oualil, Zakaria Yartaoui, Ismail Berrada \\
Afiliasi &
Mohammed VI Polytechnic University (UM6P), Maroko \\
Nama jurnal &
\textit{Expert Systems with Applications} (ESWA), Penerbit Elsevier \\
Volume / Nomor artikel &
Volume 234 / Article 121000 \\
Tahun &
2023 \\
Halaman &
Article number 121000 (artikel bernomor, bukan rentang halaman konvensional) \\
DOI &
\href{https://doi.org/10.1016/j.eswa.2023.121000}{10.1016/j.eswa.2023.121000} \\
Status Scopus &
Terindeks Scopus; ESWA merupakan jurnal \textbf{Q1} (bidang Computer Science /
Artificial Intelligence \& Engineering) dengan reputasi tinggi (CiteScore
tinggi, impact factor $>7$). \\
Kode sumber &
Tersedia publik: \url{https://github.com/meryemJanatiIdrissi/Fed-ANIDS} \\
\bottomrule
\end{tabular}
\caption{Identitas bibliografis paper Fed-ANIDS.}
\end{table}

\noindent\textbf{Alasan pemilihan.} Paper ini dipilih karena:
(1) terbit di jurnal Q1 bereputasi (ESWA, Elsevier);
(2) memadukan dua teknik unggulan sekaligus --- \emph{anomaly detection}
berbasis \emph{autoencoder} dan \emph{federated learning} --- yang menjadi
arah utama riset FL-NIDS modern;
(3) membandingkan beberapa varian \emph{autoencoder} (AE, VAE, USAD) dan dua
algoritma agregasi (FedAvg, FedProx), sehingga kaya sebagai bahan ajar;
(4) kode sumbernya terbuka sehingga mekanismenya dapat dicontohkan secara
konkret kepada mahasiswa.

%==================================================================
\section{Solusi yang Diusulkan Penulis}
%==================================================================

\subsection{Solusi}
Masalah inti yang diangkat: NIDS tradisional berbasis \emph{machine learning}
umumnya \textbf{terpusat} --- data lalu lintas jaringan dari banyak organisasi
harus dikumpulkan ke satu server untuk melatih model. Ini menimbulkan dua
persoalan: \textbf{privasi} (data lalu lintas sensitif tidak boleh keluar
organisasi) dan \textbf{ketergantungan pada data serangan berlabel} yang mahal
dan jarang.

Solusi Fed-ANIDS menggabungkan dua gagasan:
\begin{enumerate}[leftmargin=1.4em]
  \item \textbf{Federated Learning (FL)} --- setiap node melatih model
  \emph{secara lokal} dan hanya mengirim \emph{parameter model} (bukan data
  mentah) ke server agregator. Privasi data terjaga.
  \item \textbf{Anomaly Detection (AD) berbasis rekonstruksi} --- model dilatih
  \textbf{hanya pada lalu lintas normal}, lalu serangan dideteksi sebagai
  anomali ketika \emph{reconstruction error} melebihi ambang. Dengan ini,
  sistem \textbf{tidak membutuhkan data serangan berlabel} saat pelatihan.
\end{enumerate}

\subsection{Sistem}
Arsitektur sistem bersifat \emph{client--server} klasik FL:
\begin{itemize}[leftmargin=1.4em]
  \item \textbf{Klien} ($K$ node): masing-masing memiliki data lalu lintas
  lokal (hanya trafik normal untuk pelatihan). Tiap klien melatih
  \emph{autoencoder} pada datanya sendiri.
  \item \textbf{Server agregator}: menginisialisasi model global, menyiarkannya
  ke klien, lalu menggabungkan parameter yang dikembalikan klien menjadi model
  global baru. Diulang sampai konvergen.
\end{itemize}
Yang berpindah di jaringan hanyalah bobot model; data mentah tidak pernah
meninggalkan klien. Alur komunikasi antar-node ini diringkas pada
Gambar~\ref{fig:fl-arsitektur}.

\begin{figure}[h]
\centering
\begin{tikzpicture}[node distance=1.1cm and 1.9cm]
  % server
  \node[server] (srv) {\textbf{Server Agregator}\\ FedAvg / FedProx\\ model global $w^{(t)}$};
  % clients
  \node[client, below left=2.0cm and 1.6cm of srv] (c1)
        {\textbf{Klien 1}\\ data lokal\\ (trafik normal)};
  \node[client, below=2.0cm of srv] (c2)
        {\textbf{Klien 2}\\ data lokal\\ (trafik normal)};
  \node[client, below right=2.0cm and 1.6cm of srv] (c3)
        {\textbf{Klien $K$}\\ data lokal\\ (trafik normal)};
  % broadcast (down) : global weights
  \draw[flow,blue!60!black] (srv.south) -- (c1.north)
        node[midway,sloped,above,font=\scriptsize]{siar $w^{(t)}$};
  \draw[flow,blue!60!black] (srv.south) -- (c2.north);
  \draw[flow,blue!60!black] (srv.south) -- (c3.north)
        node[midway,sloped,above,font=\scriptsize]{siar $w^{(t)}$};
  % upload (up) : local weights
  \draw[flow,red!70!black] ([xshift=6mm]c1.north) -- ([xshift=-10mm]srv.south)
        node[midway,sloped,below,font=\scriptsize]{kirim $w_1^{(t)}$};
  \draw[flow,red!70!black] ([xshift=8mm]c3.north) -- ([xshift=10mm]srv.south)
        node[midway,sloped,below,font=\scriptsize]{kirim $w_K^{(t)}$};
  % privacy note
  \node[draw=none,below=0.5cm of c2,font=\scriptsize\itshape,text=gray!50!black]
        {Data mentah TIDAK pernah keluar klien --- hanya bobot model dipertukarkan.};
\end{tikzpicture}
\caption{Arsitektur \emph{federated learning} client--server. Tiap \emph{round} $t$:
server menyiarkan bobot global (panah biru, turun), klien melatih lokal lalu
mengirim bobotnya (panah merah, naik), server mengagregasi menjadi $w^{(t+1)}$.}
\label{fig:fl-arsitektur}
\end{figure}

\subsection{Model}
Fed-ANIDS mengevaluasi tiga keluarga model rekonstruksi:
\begin{itemize}[leftmargin=1.4em]
  \item \textbf{Simple Autoencoder (AE)} --- \emph{encoder}--\emph{decoder}
  yang memampatkan input ke ruang laten lalu merekonstruksinya.
  \item \textbf{Variational Autoencoder (VAE)} --- varian probabilistik dengan
  regularisasi distribusi laten.
  \item \textbf{USAD (UnSupervised Anomaly Detection)} --- dua
  \emph{autoencoder} yang dilatih secara \textbf{adversarial}; inilah varian
  yang memberi hasil terbaik pada paper ini.
\end{itemize}
Deteksi dilakukan via \textbf{intrusion score} = kesalahan rekonstruksi.
Trafik dengan skor di atas ambang $\tau$ ditandai sebagai serangan.

\subsection{Metode / Algoritma (beserta penjelasan)}

\subsubsection*{(a) Agregasi federated: FedAvg dan FedProx}
Pada tiap \emph{round} komunikasi $t$, server menyiarkan bobot global $w^{(t)}$.
Tiap klien $k$ melatih secara lokal $E$ epoch menghasilkan $w_k^{(t)}$, lalu
server menggabungkan dengan \textbf{FedAvg} (rata-rata berbobot ukuran data):
\begin{equation}
w^{(t+1)} = \sum_{k=1}^{K} \frac{n_k}{n}\, w_k^{(t)},
\qquad n = \sum_{k=1}^{K} n_k ,
\end{equation}
dengan $n_k$ jumlah sampel klien $k$. Varian \textbf{FedProx} menambahkan
\emph{proximal term} pada fungsi rugi lokal untuk menahan klien agar tidak
menyimpang terlalu jauh dari model global (penting saat data antar-klien
\emph{non-IID}):
\begin{equation}
\min_{w}\; \mathcal{L}_k(w) \;+\; \frac{\mu}{2}\,\lVert w - w^{(t)} \rVert^2 .
\end{equation}

\subsubsection*{(b) Model terbaik: USAD (adversarial autoencoder)}
USAD memakai satu \emph{encoder} $E$ bersama dan dua \emph{decoder} $D_1, D_2$
(membentuk dua \emph{autoencoder} $AE_1 = D_1\!\circ\!E$ dan
$AE_2 = D_2\!\circ\!E$). Pelatihan dua fase:
\begin{itemize}[leftmargin=1.4em]
  \item \textbf{Fase 1 (rekonstruksi).} Kedua \emph{autoencoder} belajar
  merekonstruksi input normal $x$:
  \begin{equation}
  \mathcal{L}_{AE_1} = \lVert x - AE_1(x) \rVert_2, \qquad
  \mathcal{L}_{AE_2} = \lVert x - AE_2(x) \rVert_2 .
  \end{equation}
  \item \textbf{Fase 2 (adversarial).} $AE_2$ belajar \emph{membedakan} data
  asli dari hasil rekonstruksi $AE_1$, sementara $AE_1$ berusaha
  \emph{mengelabui} $AE_2$. Permainan \emph{min--max} ini mempertajam batas
  antara normal dan anomali:
  \begin{equation}
  \min_{AE_1}\max_{AE_2}\;
  \lVert x - AE_2(AE_1(x)) \rVert_2 .
  \end{equation}
\end{itemize}
Pada inferensi, \textbf{intrusion score} menggabungkan kedua kesalahan:
\begin{equation}
s(x) = \alpha\,\lVert x - AE_1(x)\rVert_2
     + \beta\,\lVert x - AE_2(AE_1(x))\rVert_2,
\qquad \alpha+\beta = 1 .
\end{equation}
Jika $s(x) > \tau$ maka $x$ dinyatakan serangan. Ambang $\tau$ dikalibrasi dari
distribusi skor data normal (mis. persentil tinggi). Struktur encoder bersama dan
dua decoder beserta jalur adversarial digambarkan pada Gambar~\ref{fig:usad}.

\begin{figure}[h]
\centering
\begin{tikzpicture}[node distance=1.0cm and 1.4cm]
  \node[block] (x) {Input $x$};
  \node[lat, right=of x] (z) {$z$};
  \node[block, above right=0.6cm and 1.4cm of z] (d1) {$D_1$};
  \node[block, below right=0.6cm and 1.4cm of z] (d2) {$D_2$};
  \node[block, right=of d1] (w1) {$AE_1(x)$};
  \node[block, right=of d2] (w2) {$AE_2(x)$};
  \node[block, right=1.4cm of w1] (w3) {$AE_2(AE_1(x))$};
  % encoder
  \draw[flow] (x) -- (z) node[midway,above,font=\scriptsize]{Encoder $E$};
  \draw[flow] (z) -- (d1);
  \draw[flow] (z) -- (d2);
  \draw[flow] (d1) -- (w1);
  \draw[flow] (d2) -- (w2);
  % adversarial path: w1 -> E -> D2 -> w3
  \draw[flow,purple!70!black] (w1.north) to[out=60,in=120]
        node[midway,above,font=\scriptsize]{$E$ lalu $D_2$ (jalur adversarial)} (w3.north);
  % reconstruction targets
  \draw[flowb,gray!60!black] (x.south) to[out=-70,in=180]
        ([yshift=-6mm]w2.south) node[pos=0.5,below,font=\scriptsize]{rekonstruksi: $x \approx AE_1,\,AE_2$};
  \node[draw=none,below=1.3cm of z,font=\scriptsize\itshape,text=gray!50!black,align=center]
        {Encoder $E$ dipakai BERSAMA kedua decoder.\\
         Fase 1: rekonstruksi ($AE_1,AE_2$). Fase 2: adversarial ($AE_2$ vs $AE_1$).};
\end{tikzpicture}
\caption{Arsitektur USAD: satu \emph{encoder} bersama $E$ dan dua \emph{decoder}
$D_1,D_2$. Jalur ungu $AE_2(AE_1(x))$ adalah inti pelatihan adversarial yang
mempertajam batas normal--anomali.}
\label{fig:usad}
\end{figure}

\subsubsection*{(c) Alur lengkap pelatihan federated (pseudokode)}
\begin{enumerate}[leftmargin=1.6em]
  \item Server inisialisasi bobot global $w^{(0)}$ (encoder + dua decoder).
  \item \textbf{Untuk} tiap round $t = 0,1,\dots,R-1$:
  \begin{enumerate}[leftmargin=1.4em,label=\arabic*.]
    \item Server siarkan $w^{(t)}$ ke semua klien.
    \item Tiap klien $k$ latih USAD $E$ epoch \emph{pada trafik normal lokal}
    (fase rekonstruksi + adversarial) $\rightarrow w_k^{(t)}$.
    \item Klien kirim $w_k^{(t)}$ ke server (data tetap lokal).
    \item Server agregasi FedAvg/FedProx $\rightarrow w^{(t+1)}$.
  \end{enumerate}
  \item Keluaran: model global; evaluasi pada test set (normal + serangan)
  memakai \emph{intrusion score} $s(x)$ dan ambang $\tau$.
\end{enumerate}

\subsection{Desain}
Rancangan evaluasi dibuat sistematis sebagai \emph{faktorial}:
\begin{itemize}[leftmargin=1.4em]
  \item \textbf{Tiga model AD}: AE, VAE, USAD.
  \item \textbf{Dua agregasi FL}: FedAvg, FedProx.
  \item \textbf{Tiga dataset benchmark}: USTC-TFC2016, CIC-IDS2017,
  CSE-CIC-IDS2018.
  \item \textbf{Pembanding}: model terpusat (\emph{centralized}) sebagai batas
  atas, serta variasi jumlah klien untuk menguji skalabilitas.
\end{itemize}
Metrik: \emph{accuracy, precision, recall, F1-score}. Pelatihan hanya memakai
trafik normal (\emph{semi-supervised / one-class}); serangan hanya muncul saat
evaluasi. Alur deteksi saat inferensi diringkas pada Gambar~\ref{fig:deteksi}.

\begin{center}
\footnotesize
\begin{tabular}{@{}p{2.3cm}p{11.2cm}@{}}
\toprule
\textbf{Lapisan} & \textbf{Peran dalam desain} \\
\midrule
Data lokal & Trafik ternormalisasi per-klien; hanya kelas normal untuk latih. \\
Model lokal & USAD (encoder bersama + 2 decoder), ringan agar hemat komunikasi. \\
Agregasi & FedAvg (dasar) / FedProx (stabil saat non-IID). \\
Deteksi & Intrusion score dari reconstruction error + ambang $\tau$. \\
\bottomrule
\end{tabular}
\end{center}

\begin{figure}[h]
\centering
\begin{tikzpicture}[node distance=0.9cm and 1.5cm]
  \node[block,fill=green!10] (traf) {Trafik masuk $x$};
  \node[block, right=of traf] (model) {Model USAD\\ global};
  \node[block, right=of model] (err) {Hitung\\ \emph{recon. error}};
  \node[block, right=of err] (score) {Intrusion\\ score $s(x)$};
  \node[diamond,draw=black,fill=yellow!20,aspect=2,font=\scriptsize,
        right=1.3cm of score] (dec) {$s(x) > \tau$?};
  \node[block,fill=red!15, above right=0.4cm and 1.3cm of dec] (atk) {SERANGAN};
  \node[block,fill=green!15, below right=0.4cm and 1.3cm of dec] (norm) {NORMAL};
  \draw[flow] (traf) -- (model);
  \draw[flow] (model) -- (err);
  \draw[flow] (err) -- (score);
  \draw[flow] (score) -- (dec);
  \draw[flow] (dec) -- (atk) node[midway,sloped,above,font=\scriptsize]{ya};
  \draw[flow] (dec) -- (norm) node[midway,sloped,below,font=\scriptsize]{tidak};
\end{tikzpicture}
\caption{Alur deteksi (inferensi): trafik direkonstruksi model global, kesalahan
rekonstruksi diubah jadi \emph{intrusion score} $s(x)$, lalu dibandingkan dengan
ambang $\tau$. Model dilatih hanya pada trafik normal, sehingga serangan muncul
sebagai anomali berskor tinggi.}
\label{fig:deteksi}
\end{figure}

\subsection{Keunggulan Sistem}
\begin{enumerate}[leftmargin=1.4em]
  \item \textbf{Menjaga privasi} --- data mentah tak pernah keluar klien; hanya
  bobot yang dipertukarkan. Cocok untuk kolaborasi antar-organisasi.
  \item \textbf{Tidak butuh label serangan} --- pelatihan \emph{one-class} pada
  trafik normal; mampu mendeteksi serangan baru (\emph{zero-day}) sebagai
  anomali.
  \item \textbf{Akurasi tinggi \& stabil} --- varian USAD memberi F1 mendekati
  model terpusat, dengan deteksi tajam berkat pelatihan adversarial.
  \item \textbf{Efisien komunikasi} --- \emph{autoencoder} ringan; FedProx
  mempercepat kestabilan konvergensi pada data non-IID.
  \item \textbf{Teruji lintas-dataset} --- konsisten di tiga benchmark NIDS
  yang berbeda karakteristik.
\end{enumerate}

%==================================================================
\section{Contoh Kode: Inti Metode Terbaik (USAD + FedAvg)}
%==================================================================

Berikut ilustrasi \emph{inti} metode (disederhanakan untuk pembelajaran, PyTorch
dan NumPy). Kode difokuskan pada tiga bagian kunci: (a) model USAD, (b)
pelatihan lokal dua fase, dan (c) agregasi FedAvg. Ini bukan salinan verbatim
repositori penulis, melainkan rekonstruksi edukatif agar mahasiswa memahami
mekanismenya.

\subsection{(a) Arsitektur USAD}
\begin{lstlisting}[style=py,caption={Encoder bersama + dua decoder (USAD).}]
import torch, torch.nn as nn

class Encoder(nn.Module):
    def __init__(self, d_in, d_lat):
        super().__init__()
        self.net = nn.Sequential(
            nn.Linear(d_in, d_in // 2), nn.ReLU(),
            nn.Linear(d_in // 2, d_lat), nn.ReLU())
    def forward(self, x): return self.net(x)

class Decoder(nn.Module):
    def __init__(self, d_lat, d_out):
        super().__init__()
        self.net = nn.Sequential(
            nn.Linear(d_lat, d_out // 2), nn.ReLU(),
            nn.Linear(d_out // 2, d_out), nn.Sigmoid())
    def forward(self, z): return self.net(z)

class USAD(nn.Module):
    """Satu encoder E dipakai bersama; dua decoder D1, D2."""
    def __init__(self, d_in, d_lat):
        super().__init__()
        self.E  = Encoder(d_in, d_lat)
        self.D1 = Decoder(d_lat, d_in)
        self.D2 = Decoder(d_lat, d_in)
    def forward(self, x):
        z   = self.E(x)
        w1  = self.D1(z)            # AE1(x)
        w2  = self.D2(z)            # AE2(x)
        w3  = self.D2(self.E(w1))   # AE2(AE1(x))
        return w1, w2, w3
\end{lstlisting}

\subsection{(b) Pelatihan lokal dua fase (rekonstruksi + adversarial)}
\begin{lstlisting}[style=py,caption={Satu klien melatih USAD pada trafik normal lokal.}]
def train_local(model, X_normal, epochs=5, lr=1e-3):
    opt = torch.optim.Adam(model.parameters(), lr=lr)
    for ep in range(1, epochs + 1):
        n = ep  # bobot adversarial meningkat tiap epoch
        for xb in X_normal:                       # xb: batch trafik normal
            w1, w2, w3 = model(xb)
            # Fase 1: rekonstruksi. Fase 2: adversarial (min-max).
            loss_ae1 = (1/n)   * torch.mean((xb - w1)**2) \
                     + (1-1/n) * torch.mean((xb - w3)**2)
            loss_ae2 = (1/n)   * torch.mean((xb - w2)**2) \
                     - (1-1/n) * torch.mean((xb - w3)**2)
            loss = loss_ae1 + loss_ae2
            opt.zero_grad(); loss.backward(); opt.step()
    return model.state_dict()                     # hanya bobot yang dikirim

def intrusion_score(model, x, alpha=0.5, beta=0.5):
    w1, _, w3 = model(x)
    return alpha*torch.mean((x-w1)**2, dim=1) \
         + beta *torch.mean((x-w3)**2, dim=1)      # s(x) > tau -> serangan
\end{lstlisting}

\subsection{(c) Agregasi FedAvg di server}
\begin{lstlisting}[style=py,caption={Rata-rata berbobot ukuran data (FedAvg).}]
import copy, torch

def fedavg(client_states, client_sizes):
    total = float(sum(client_sizes))
    avg = copy.deepcopy(client_states[0])
    for key in avg.keys():
        avg[key] = sum((client_sizes[k]/total) * client_states[k][key]
                       for k in range(len(client_states)))
    return avg                                     # -> bobot global baru

def federated_round(global_model, clients, rounds=50, epochs=5):
    for t in range(rounds):
        states, sizes = [], []
        for Xk in clients:                         # tiap klien = trafik normal lokal
            local = copy.deepcopy(global_model)
            local.load_state_dict(global_model.state_dict())  # terima bobot global
            states.append(train_local(local, Xk, epochs))
            sizes.append(len(Xk))
        global_model.load_state_dict(fedavg(states, sizes))   # agregasi
    return global_model
\end{lstlisting}

\noindent\textbf{Catatan pengajaran.} Perhatikan: fungsi \texttt{train\_local}
mengembalikan \emph{state\_dict} (bobot), \textbf{bukan} data. Inilah inti
privasi FL. Faktor \texttt{1/n} vs \texttt{1-1/n} menyeimbangkan tujuan
rekonstruksi (awal pelatihan) dan adversarial (akhir pelatihan) --- ciri khas
USAD.

%==================================================================
\section{Hasil Percobaan yang Dilakukan}
%==================================================================

Ringkasan temuan utama paper (dari evaluasi pada tiga dataset benchmark):

\begin{itemize}[leftmargin=1.4em]
  \item \textbf{USAD unggul.} Di antara ketiga model AD, varian berbasis
  \textbf{USAD} secara konsisten memberi \emph{F1-score} tertinggi, mengungguli
  AE dan VAE pada ketiga dataset.
  \item \textbf{Mendekati model terpusat.} Model \emph{federated} (USAD) mencapai
  kinerja yang \textbf{sangat dekat} dengan model terpusat (\emph{centralized})
  --- selisih kecil --- padahal data \textbf{tidak pernah dipusatkan}. Ini bukti
  bahwa privasi dapat dijaga nyaris tanpa mengorbankan akurasi.
  \item \textbf{Skor tinggi lintas-dataset.} Pada CIC-IDS2017, CSE-CIC-IDS2018,
  dan USTC-TFC2016, metrik \emph{accuracy}, \emph{precision}, \emph{recall}, dan
  \emph{F1} berada pada kisaran tinggi (umumnya di atas $\sim$0,95 untuk
  konfigurasi terbaik), menunjukkan generalisasi yang baik.
  \item \textbf{FedProx vs FedAvg.} FedProx membantu kestabilan ketika distribusi
  antar-klien tidak seragam (\emph{non-IID}), meski pada kondisi relatif seimbang
  perbedaannya tipis terhadap FedAvg.
  \item \textbf{Robust terhadap serangan tak-terlihat.} Karena dilatih hanya pada
  trafik normal, sistem mendeteksi beragam jenis serangan sebagai anomali tanpa
  perlu contoh serangan saat pelatihan.
\end{itemize}

\begin{table}[h]
\centering
\footnotesize
\renewcommand{\arraystretch}{1.25}
\begin{tabular}{@{}lcccc@{}}
\toprule
\textbf{Konfigurasi (ilustratif)} & \textbf{Accuracy} & \textbf{Precision} &
\textbf{Recall} & \textbf{F1} \\
\midrule
Centralized (USAD)        & tinggi & tinggi & tinggi & \textbf{tertinggi} \\
Federated USAD (FedAvg)   & $\approx$ centralized & tinggi & tinggi & $\approx$ centralized \\
Federated USAD (FedProx)  & $\approx$ centralized & tinggi & tinggi & stabil (non-IID) \\
Federated AE / VAE        & baik & baik & baik & di bawah USAD \\
\bottomrule
\end{tabular}
\caption{Pola hasil (deskriptif). Angka pasti lihat paper asli; tabel ini
menegaskan \emph{tren}: USAD terbaik, federated $\approx$ centralized.}
\end{table}

\noindent\textit{Catatan kejujuran akademik:} nilai metrik persisnya mengacu
pada tabel di paper asli (DOI di atas). Tabel ini sengaja menyajikan \emph{pola}
agar fokus pembelajaran pada mekanisme, bukan menghafal angka.

%==================================================================
\section{Kesimpulan}
%==================================================================

\begin{enumerate}[leftmargin=1.4em]
  \item \textbf{Fed-ANIDS} menunjukkan bahwa menggabungkan \emph{federated
  learning} dengan \emph{anomaly detection} berbasis \emph{autoencoder}
  merupakan salah satu pendekatan paling unggul untuk NIDS yang
  \textbf{menjaga privasi} sekaligus \textbf{akurat}.
  \item Model \textbf{USAD} (autoencoder adversarial) memberi kinerja terbaik;
  pelatihan dua fase (rekonstruksi lalu adversarial) menghasilkan batas
  normal--anomali yang tajam.
  \item Kinerja \emph{federated} \textbf{mendekati} model terpusat tanpa pernah
  memusatkan data mentah --- bukti nyata bahwa privasi dan akurasi tidak harus
  dikorbankan satu sama lain.
  \item Pendekatan \emph{semi-supervised} (hanya butuh trafik normal) membuat
  sistem praktis di dunia nyata, di mana data serangan berlabel langka, dan
  mampu menangkap serangan baru sebagai anomali.
  \item \textbf{Pelajaran untuk FL-NIDS secara umum:} kunci keberhasilan ada
  pada (a) pemilihan model lokal yang tepat (USAD $>$ VAE $>$ AE), (b) algoritma
  agregasi yang tahan \emph{non-IID} (FedProx), dan (c) strategi deteksi berbasis
  rekonstruksi yang tidak bergantung label serangan.
\end{enumerate}

\vspace{1em}
\hrule
\vspace{0.5em}
\noindent\footnotesize
\textbf{Referensi utama:} M. Janati Idrissi dkk., ``Fed-ANIDS: Federated
learning for anomaly-based network intrusion detection systems,''
\textit{Expert Systems with Applications}, vol.~234, art.~121000, 2023.
DOI: \href{https://doi.org/10.1016/j.eswa.2023.121000}{10.1016/j.eswa.2023.121000}.
Metode USAD mengacu pada J. Audibert dkk., ``USAD: UnSupervised Anomaly
Detection on Multivariate Time Series,'' \textit{ACM SIGKDD}, 2020.

\end{document}
